% CAST 논문 수식 소스 — Word 붙여넣기용
% 마지막 갱신: 2026-09-16
%
% 쓰는 법
%   Word 2016 이상: 삽입 → 수식 → (수식 도구) LaTeX 선택 → 아래 한 줄을 붙여넣고 변환
%   각 식의 $...$ 안쪽만 복사한다. % 로 시작하는 줄은 설명이니 복사하지 않는다.
%
% 왜 LaTeX 인가
%   현재 CAST.docx 의 수식 12개가 전부 일반 텍스트다 (Word 수식 개체 0개).
%   그래서 수식처럼 보이지 않는다. LaTeX 로 넣으면 학술지 조판과 같은 모양이 된다.
%
% 주의
%   Word 의 LaTeX 모드는 일부 명령을 지원하지 않는다.
%   \mathit, \text, \operatorname 등이 안 먹으면 아래 [Word 대안] 을 쓴다.
%   변환 후에는 반드시 눈으로 확인할 것.
%
% 기호 규약은 공식.txt [0] 을 따른다. 특히
%   미배정 변수는 z_j  (기존 s_j 는 제출 시각과 충돌하므로 바꿨다)
%   예측은 발행 시각을 윗첨자로 명시: \hat{C}_r^{(\tau)}(t)


% ══════════════════════════════════════════════════════════════
% 1. 이미 논문에 있는 식 — 텍스트를 수식 개체로 바꾸기만 하면 된다
% ══════════════════════════════════════════════════════════════

% (1) 목적함수 — 총 배출량 최소화                      문단 446
$\min \sum_{j \in J} P \int_{\tau_j}^{\tau_j + d_j} C_{r(j)}(t)\, dt$

% (2) 모든 작업은 정확히 한 리전에 배정                 문단 448
$\sum_{r \in R} x_{jr} = 1, \quad \forall j \in J$

% (3) 용량 제약 — 이 식의 \forall t 가 연속 시각이다     문단 449
%     2026-09-16 구현 버그를 판정한 근거가 이 식이었다
$\left| \{\, j : r(j) = r,\ \tau_j \le t < \tau_j + d_j \,\} \right| \le \left\lfloor \eta \cdot cap_r \right\rfloor, \quad \forall r \in R,\ \forall t$

% (4) 실행은 제출 이후이고 마감 이내                    문단 450
$s_j \le \tau_j, \quad \tau_j + d_j \le D_j, \quad \forall j \in J$

% (5) 시간 지연과 네트워크 지연의 합이 예산 이내         문단 451
%     본문이 밝힌 대로 구현에서 강제되지 않고 수치상 구속되지도 않는다
$(\tau_j - s_j) + \ell_{o_j r(j)} \le L_j, \quad \forall j \in J$

% (6) 정규화 — 탄소는 슬롯별, 지연은 고정 분모          문단 475
$\tilde{M}_r = \frac{\hat{C}_r}{\max_{r'} \hat{C}_{r'}}, \qquad \tilde{\ell}_{or} = \frac{\ell_{or}}{244}$

% (7) 슬롯 ILP 목적함수                                문단 479
$\min \sum_{j} \sum_{r} \left( \alpha \tilde{M}_r + (1 - \alpha) \tilde{\ell}_{o_j r} \right) x_{jr} + P_{slack} \sum_{j} z_j$

% (8) 슬롯 ILP 제약                                    문단 480
$\sum_{r} x_{jr} + z_j = 1 \quad \forall j, \qquad \sum_{j} x_{jr} \le a_r \quad \forall r$

% (9) 파레토 무릎점 — 기여 1                            문단 484
$\alpha^{*} = \arg\min_{\alpha} \sqrt{ \tilde{\lambda}(\alpha)^{2} + \tilde{\mu}(\alpha)^{2} }$

% (10) 탄소 비용과 지연 비용의 0-1 정규화               문단 499
$\hat{C}(t) = \frac{\bar{C}(t) - C_{\min}}{C_{\max} - C_{\min}}, \qquad \hat{D}(t) = \frac{t - t_{early}}{t_{late} - t_{early}}$

% (11) 슬롯 점수 — 시간 이동의 판단 기준                문단 500
$score(t) = \alpha(k)\, \hat{C}(t) + \bigl(1 - \alpha(k)\bigr) \hat{D}(t), \qquad \alpha(k) = \frac{6 - k}{5}$

% (12) 탄소 회계 — 배출량은 실측값으로 산정             문단 514
%      ** 이미 있는 식이다. 새 식은 (13) 부터다. **
$E_j = P \int_{\tau_j}^{\tau_j + d_j} C_{r(j)}(t)\, dt$


% ══════════════════════════════════════════════════════════════
% 2. 새 식 — 용량 인지 온라인 시간 이동 (§3.6 신설)
% ══════════════════════════════════════════════════════════════

% (13) 가용 자리 — 식 (3) 을 가용 자리 형태로 다시 쓴 것
%      논문에 이렇게 쓸 것:
%      "5.2절에서 식 (3) 으로 선언했으나 구현하지 않았던 제약을
%       이 알고리즘이 실제로 강제한다."
$n_r(t) = \left| \{\, j : r(j) = r,\ \tau_j \le t < \tau_j + d_j \,\} \right|, \qquad a_r(t) = \left\lfloor \eta \cdot cap_r \right\rfloor - n_r(t)$

% (14) 희망 집합 — 남은 윈도우에서 지금이 최선인 작업
%      score 는 식 (11) 을 그대로 쓴다. 새 점수 함수를 만들지 않는다.
%      단 예측은 발행 시각이 t 인 것을 쓴다 (규약 R4).
$Q(t) = \{\, j : s_j \le t,\ \tau_j \text{ 미정} \,\}$

$P_r(t) = \Bigl\{\, j \in Q(t) : r(j) = r,\ \arg\min_{t' \in W_j(t)} score_j(t') = \max(t,\, s_j) \,\Bigr\}$

% (15) 우선순위와 선택 — 순수 EDF
%      탄소 이득으로 정렬하면 나중에 측정할 지표로 승자를 미리 뽑는 셈이라
%      순환 논리가 된다. 여유만 보면 그 문제가 없고, 여유 0 이 자동으로 1순위다.
$u_j(t) = (D_j - d_j) - t$

$F_r(t) = \{\, j \in Q(t) : r(j) = r,\ u_j(t) < 1 \,\}$

$S_r(t) = F_r(t) \;\cup\; R_r(t)\bigl[\,1 \ldots \max\bigl(a_r(t) - |F_r(t)|,\ 0\bigr)\,\bigr]$

% R_r(t) 는 P_r(t) \ F_r(t) 를 u_j(t) 오름차순으로 정렬한 것


% ══════════════════════════════════════════════════════════════
% 3. 규약 — 번호는 안 붙이되 본문에 반드시 적을 것
%    2026-09-16 구현 버그가 전부 여기서 나왔다 (공식.txt [2] 참고)
% ══════════════════════════════════════════════════════════════

% R1 후보 슬롯 격자 — 벽시계 정시에 현재 시점 자체를 더한다
%    첫 항이 "지금 당장 실행" 이다. 이게 없으면 제출 시각이 정시가 아닌 작업이
%    다음 정시까지 밀리고, 지연 예산이 1시간 미만인 작업은 마감을 어긴다.
$W_j(t) = \bigl\{ \max(t,\, s_j) \bigr\} \cup \bigl\{\, t' \in \mathbb{Z} : \lceil \max(t,\, s_j) \rceil \le t' \le \min(D_j - d_j,\ t + H) \,\bigr\}$

% R4 예측 발행 시각 — 두 알고리즘의 차이가 정확히 여기다
%    사전 예약형은 제출 시점 예측 하나로 전 윈도우를 평가하고
%    롤링형은 평가하는 매 슬롯의 최신 예측으로 다시 평가한다
$\bar{C}_j^{(\tau)}(t') = \frac{1}{d_j} \int_{t'}^{t' + d_j} \hat{C}_{r(j)}^{(\tau)}(u)\, du$


% ══════════════════════════════════════════════════════════════
% 4. Word 대안 — LaTeX 모드가 명령을 못 알아들을 때
% ══════════════════════════════════════════════════════════════
%
% Word 수식 편집기는 UnicodeMath(선형 서식)도 받는다.
% 수식 안에서 아래처럼 치고 스페이스를 누르면 변환된다.
%
%   \lfloor  ->  ⌊        \rfloor  ->  ⌋
%   \tau     ->  τ        \alpha   ->  α
%   \hat{C}  ->  C 위에 모자   (C\hat 다음 스페이스)
%   \sum_    ->  ∑ 아래첨자
%   \le      ->  ≤        \forall  ->  ∀
%   a_r(t)   ->  아래첨자 r        (a_r 다음 스페이스)
%
% 지금 본문이 쓰는 τ_(j) 같은 표기가 이미 UnicodeMath 와 비슷하다.
% 수식 개체 안에 넣고 스페이스만 누르면 상당수가 그대로 변환된다.
%
%
% ══════════════════════════════════════════════════════════════
% 5. 남은 작업
% ══════════════════════════════════════════════════════════════
%
% [ ] 기존 식 (1)~(12) 를 전부 Word 수식 개체로 교체   ← 지금 0개다
% [ ] 식 (7)(8) 의 s_j 를 z_j 로 (제출 시각과 충돌)
% [ ] 식 (8) 의 avail_r 표기를 a_r 로 통일
% [ ] 새 식 (13)(14)(15) 를 §3.6 에 삽입
% [ ] R1 · R4 를 본문 문장으로 명시
% [ ] 본문에서 식 번호를 산문으로 인용 (지금 문단 452·454·456 에서만 부른다.
%     식 (6)~(12) 는 한 번도 번호로 참조되지 않는다.
%     Caspian·CarbonFlex 는 알고리즘 블록 밖 산문에서 식 번호를 인용한다)
